\documentclass[11pt,a4paper,fontset=fandol]{ctexart}
\usepackage[margin=2.2cm]{geometry}
\usepackage{amsmath,amssymb}
\usepackage{booktabs,longtable,tabularx,array}
\usepackage{enumitem}
\usepackage{xcolor}
\usepackage{hyperref}
\usepackage{xurl}
\usepackage{microtype}

\hypersetup{colorlinks=true,linkcolor=blue!55!black,urlcolor=blue!60!black,citecolor=blue!55!black}
\setlist{nosep,leftmargin=2em}
\setlength{\parindent}{2em}
\setlength{\parskip}{0.35em}
\newcommand{\code}[1]{\texttt{#1}}
\newcommand{\pzero}{\textbf{P0}}
\newcommand{\pone}{\textbf{P1}}
\newcommand{\ptwo}{\textbf{P2}}

\title{从静态基准到可进化环境：\\基于 EnvHarness 的环境侧演化设计与最小落地方案}
\author{Env-Agent-RSI 项目研究报告}
\date{2026年9月26日}

\begin{document}
\maketitle

\begin{abstract}
本文把研究范围严格限定在环境侧：不修改模型权重、Agent 提示词、记忆或技能，而是改变 Agent 所面对的初始状态、动作与观察契约、转移故障、资源约束和任务组合。EnvHarness 的核心价值不是“再生成一批任务”，而是把已有、可验证、可重置的静态环境变成可组合的环境变换层。本文首先解释 EnvHarness 的 Stage/Contract/Chain 抽象及其开源实现中的 Setup/Rules/Link 映射；随后比较程序化随机化、课程选择、regret 驱动环境设计、进化式关卡编辑和 LLM 生成可执行环境等路线；最后提出本项目的首版策略：使用受类型约束的环境变换算子、独立验证器、确定性快照和非支配候选档案，先在一个本地、状态化 API 微环境中跑通环境生成、验证、回放和谱系记录，再接入 AppWorld、$\tau$-bench 或 ALFWorld。配套 Excel 工作簿给出相关 benchmark 的例子、优先级、评分及来源。
\end{abstract}

\section{研究边界与结论}

本阶段回答的问题是：\textbf{环境如何自动产生有用、可解、可复现、可验证的变化？} 这里的“环境进化”不等同于 Agent 自我进化。环境生成器可以读取 Agent 的轨迹来估计能力边界，但候选修改只能落在环境规范上，且 Agent 本身在同一轮比较中保持冻结。

建议的第一版结论如下：

\begin{enumerate}
  \item 不直接在开放文本空间里用 MCTS 搜索“任意新环境代码”。首版先定义一个小而闭合的\textbf{类型化变换语言}，对初始状态、动作、观察、转移故障和预算做局部修改。
  \item 搜索结构使用\textbf{带谱系的候选档案（DAG）}，而不是强制二叉树。一个候选可以由多个父环境的变换组合而成；相同环境规范通过哈希去重。
  \item 每个候选必须依次通过：可执行、状态隔离、验证器不变、Oracle 可解、故障确实触发、重复回放一致、难度处于目标区间。
  \item 首个环境选择状态化 API 微环境，围绕“重复调用、提交前失败、提交后响应丢失、陈旧读取、幂等键”建立 8--12 个变体。它比浏览器或桌面环境更便宜，也最贴合当前关注的工具调用问题。
  \item Stage 与 Contract 是首版主线；Chain 延后。EnvHarness 论文自己的自动 EnvRigger 实验也排除了 Chain，因为跨环境内部状态难以观测和验证\cite{envharness}。
\end{enumerate}

\section{EnvHarness 具体改变了什么}

EnvHarness 将环境写成
\begin{equation}
E=(\mathcal{S},\mathcal{A},\mathcal{O},T,R,s_0),
\end{equation}
其中 $\mathcal{S}$ 是状态空间，$\mathcal{A}$ 是动作空间，$\mathcal{O}$ 是观察空间，$T$ 是转移函数，$R$ 来自原生验证器，$s_0$ 是初始状态。环境组件 $w$ 在标准接口边界产生新环境 $E'=w(E)$，而不直接修改底层模拟器\cite{envharness}。

\subsection{Stage / Setup：通过合法动作重放改变起点}

Stage 用动作序列 $\delta=(a_1,\ldots,a_k)$ 在重置后推进环境：
\begin{equation}
s'_0=T(\cdots T(T(s_0,a_1),a_2)\cdots,a_k).
\end{equation}
这有两个重要性质。第一，新起点是\textbf{可达状态}，不是任意篡改内部对象；第二，同一动作序列可回放，因此验证阶段能够检查确定性。它既可增加难度，例如把目标对象藏进关闭的抽屉，也可预先完成前置子目标来缩短任务。

\subsection{Contract / Rules：重写交互契约}

Contract 对动作、观察和转移分别施加 $f_A,f_O,f_T$：
\begin{equation}
(\mathcal{A}',\mathcal{O}',T')=(f_A(\mathcal{A}),f_O(\mathcal{O}),f_T(T)),
\end{equation}
并保留原生终局验证器 $R$。典型用法包括：屏蔽捷径动作、要求前置条件、截断观察、注入超时、把“写入已成功但响应丢失”作为后置故障，以及限制剩余步数。对本项目而言，重复调用问题主要属于 $f_T$：同样是“工具返回错误”，必须进一步区分\textbf{提交前失败}和\textbf{提交后失败}，否则无法训练幂等恢复策略。

\subsection{Chain / Link：组合更长任务}

Chain 用组合逻辑 $g$ 把 $E$ 与 $E_{\mathrm{ext}}$ 连接成 $E'=g(E,E_{\mathrm{ext}})$。串行 Chain 的复合成功条件可由各子任务验证器的合取构造。论文也指出，分支或交错组合若共享中间状态，就不能仅靠两个终局 verdict 合成可信验证器；还需要语义兼容性和新的复合目标验证器\cite{envharness}。因此首版不应把 Chain 当作自动生成主线。

\subsection{论文术语与代码术语}

\begin{center}
\begin{tabular}{lll}
\toprule
论文术语 & 当前开源代码 & 生成字段 \\
\midrule
Stage & Setup / Setups & \code{in\_env\_actions} \\
Contract & Rules & \code{rules\_code} \\
Chain & Link & 暂无 EnvRigger 直接生成字段 \\
\bottomrule
\end{tabular}
\end{center}

开源实现以统一的 \code{ActionableEnv} 接口适配 ALFWorld、WebArena、SWE-bench Verified、OfficeQA 与 SpreadsheetBench，并要求 \code{reset/step/observe/evaluate/get\_env\_state/save\_state/from\_state}\cite{envharnessrepo}。组件只能读取可序列化状态视图，不能直接持有浏览器、容器或 socket 句柄。这一限制正是跨 benchmark 复用的基础。

\section{环境可进化的操作空间}

\begin{longtable}{p{2.5cm}p{4.4cm}p{4.2cm}p{3.5cm}}
\toprule
轴 & 可修改内容 & 示例 & 必须保持的约束 \\
\midrule
\endfirsthead
\toprule
轴 & 可修改内容 & 示例 & 必须保持的约束 \\
\midrule
\endhead
初始状态 & 通过动作重放或声明式 seed 改变起点 & 已有一条订单、对象位于关闭容器、第一子目标已完成 & 状态可达；重置可复现；不泄漏答案 \\
动作契约 & 过滤、重命名、参数约束、前置条件 & 禁止直接 URL；写操作必须带幂等键 & 被阻止动作有明确语义；不能破坏 Oracle 路径 \\
观察契约 & 屏蔽、压缩、分页、延迟、噪声 & 只返回前两段；读取返回陈旧版本 & 关键信息仍可通过合法交互获得 \\
转移故障 & 在调用前后注入失败或副作用 & 提交前超时、提交后响应丢失、部分写入 & 故障时点可记录；状态与返回必须一致 \\
资源与预算 & 步数、时间、并发、查询和 token 限额 & 全量测试超时，只允许目标测试 & 预算应可审计且不改变终局目标 \\
任务数据 & 改参数、实体、目标和背景记录 & 不同用户、订单、日期、对象数量 & 参数合法；验证器覆盖新实例 \\
组合 & 串行、分支或交错子任务 & 更新记录后发送通知 & 复合 verifier 完整；共享状态语义明确 \\
环境程序 & 生成工具、数据库、网页或模拟器代码 & 合成新 API 服务或网页 & 可执行、可审计、可解、无验证漂移 \\
课程选择 & 从已有或变异候选中选择下一环境 & 选择当前成功率 20--80\% 的候选 & 选择器不应把 Agent 改动混入环境对照 \\
\bottomrule
\end{longtable}

\section{其他环境设计路线及其启示}

\subsection{程序化生成与域随机化}

Procgen 提供 16 个程序化游戏，用大量关卡测试样本效率和未见关卡泛化\cite{procgen}。这类方法的优势是吞吐量高、生成边界清楚；缺点是变化空间由人工提前编码，且均匀随机不一定命中当前 Agent 的薄弱处。对本项目的启示是：\textbf{先把合法变换空间写清楚，再谈智能搜索}。

\subsection{重放与课程选择}

Prioritized Level Replay（PLR）根据学习潜力重采样已有程序化关卡，而不是生成全新代码\cite{plr}。它对应本项目最便宜的控制组：先维护环境档案，再根据失败率、学习进展或覆盖度选择下一批环境。很多情况下，选择比生成更稳定。

\subsection{Regret 驱动的环境生成}

PAIRED 用 protagonist 与 antagonist 的表现差估计 regret，让教师产生“困难但仍可解”的环境\cite{paired}。ACCEL 不从空白生成，而是持续编辑已有可行关卡，并用 regret 把课程保持在能力前沿\cite{accel}。这两项工作支持一个关键设计：候选分数不应只最大化失败率；应最大化\textbf{能力差距，同时保留可解性}。

\subsection{开放式共同进化}

POET 同时维护环境与配对策略的种群，并允许策略在环境之间迁移；Enhanced POET 用 novelty 和“最终被解决”的条件记录真正的新挑战\cite{poet}。它的价值在于保留多条进化路径而非贪心保留单一最优节点，但计算代价很高。首版可借用其“档案 + 谱系 + 非支配保留”思想，不复制其大规模共同进化。

\subsection{LLM 生成环境}

EnvGen 让 LLM 生成和调整具身环境配置\cite{envgen}；GenEnv 用生成式模拟器围绕近端发展区产生任务\cite{genenv}；Agent-World 从工具生态、数据库和主题合成可执行任务与环境\cite{agentworld}；VeriEnv 克隆网站并用后端状态提供程序化验证\cite{verienv}；OMNI-EPIC 直接生成环境与奖励代码\cite{omniepic}。这些路线扩展性强，但都把可执行性、真实性、可解性与 verifier 正确性变成额外工程问题。EnvHarness 的相对优势在于复用原生 verifier，代价是变化受现有环境能力边界约束。

\section{建议的环境搜索策略}

\subsection{不用开放式树搜索，使用受约束候选档案}

棋类 MCTS 的动作集有限，Agent 环境代码没有天然边界。解决办法不是强行把每个自然语言想法当作 action，而是先定义有限的算子族：

\begin{equation}
\mathcal{M}=\{\texttt{setup.replay},\texttt{action.block},\texttt{action.require},
\texttt{observation.mask},\texttt{transition.fail},\texttt{budget.limit}\}.
\end{equation}

每个算子有 JSON Schema、前置条件、影响范围和逆变换。一个环境候选是基础环境加有序变换列表。候选间保留父子关系，但数据结构采用 DAG：组合候选可以有多个父节点，规范哈希相同的候选合并为一个节点。

\subsection{生成循环}

\begin{enumerate}
  \item \textbf{收集边界证据：}在冻结 Agent 上运行基线环境，聚合成功、失败、循环、重复调用、无效参数、提前终止和步数浪费。
  \item \textbf{归因：}把失败映射到预定义环境轴，例如“提交后响应丢失导致重复写入”映射到 \code{transition.fail(after\_commit)}。
  \item \textbf{提案：}从一个已验证父候选出发，只允许一次加入 1--2 个类型化变换；复杂候选通过逐步组合产生。
  \item \textbf{静态门禁：}检查 schema、作用对象、禁止调用、资源上限及 verifier 独立性。
  \item \textbf{动态门禁：}运行 Oracle、冻结 Agent 与随机/弱基线；检查重置、快照、故障触发和状态隔离。
  \item \textbf{评分与归档：}只保留非支配候选，记录规范、父节点、轨迹、指标、验证结果和生成成本。
\end{enumerate}

\subsection{评分函数}

建议不要只用成功率，而使用多目标向量：
\begin{align}
q(E') = (&- |\hat p_{\mathrm{agent}}-p^*|,\quad
\hat p_{\mathrm{oracle}},\quad
\mathrm{trigger\_rate},\quad
\mathrm{novelty},\\
&-\mathrm{runtime\_cost},\quad
-\mathrm{nondeterminism},\quad
-\mathrm{collateral\_change}).
\end{align}
其中 $p^*$ 初期可设为 0.4--0.7，避免候选过易或完全不可解。候选先通过硬门禁，再进入 Pareto 档案。若必须选择单个扩展父节点，再对已归一化指标加权，而不是在验证前混成一个分数。

\subsection{为什么不是二叉问题树}

“工具调用错误/工具返回错误”适合做\textbf{诊断 taxonomy}，不适合直接充当环境搜索树的唯一结构。一个真实候选往往同时涉及故障时点、可观察性、幂等性和预算；它也可能同时修复多个父类问题。二叉树会把互相正交的轴强行排成层级，并重复生成等价叶子。更合适的是：taxonomy 用于标注失败，DAG 用于记录环境规范的演化，benchmark task 用于独立验证。

\section{Benchmark 选择原则与优先级}

配套工作簿 \code{benchmark\_landscape.xlsx} 按六项各 0--2 分评分：状态可控性、验证器质量、重置能力、环境变换适配度、搭建成本、与项目目标的贴合度。总分 10--12 为 P0，7--9 为 P1，0--6 为 P2。评分是本项目启动决策，不代表 benchmark 的学术质量。

\begin{itemize}
  \item \pzero：本地状态化 API 微环境、EnvHarness toy24 smoke、AppWorld、$\tau$-bench。共同特点是状态与工具语义清晰，便于注入提交前/后故障并做终态验证。AppWorld 有 9 个应用、457 个 API 和状态单测式 evaluator\cite{appworld}；$\tau$-bench 直接比较对话后的数据库状态\cite{taubench}。
  \item \pone：ALFWorld、WebShop、WebArena/BrowserGym、SWE-bench Verified、Terminal-Bench、OfficeQA、SpreadsheetBench、AgentBench/AgentBoard。它们用于验证跨域性，但安装或运行成本更高，或对环境变换的支持需要额外 Bridge。
  \item \ptwo：OSWorld、AndroidWorld、Pokemon、Factorio、开放式物理环境等。它们很真实，但不适合最初验证环境包装逻辑；应在状态快照、隔离和成本控制成熟后接入。OSWorld 和 AndroidWorld 的任务均带程序化初始化与执行验证，但需要完整 GUI/设备运行时\cite{osworld,androidworld}。
\end{itemize}

Agent RSI 工作常用的 benchmark 也列在工作簿中，目的是理解评价习惯而非首版接入。例如 DGM 使用 SWE-bench 与 Polyglot\cite{dgm}，SICA 使用 SWE-bench Verified、LiveCodeBench 和合成 agent benchmark\cite{sica}，AFlow 使用 HotpotQA、DROP、HumanEval、MBPP、GSM8K 与 MATH\cite{aflow}。这些工作主要搜索 Agent 或 workflow，本项目暂时只把它们作为对照。

\section{Minimal 可跑通环境计划}

\subsection{首个任务：Exactly-once Append}

环境维护一个本地 JSON 或 SQLite 状态：\code{items(id, value, idempotency\_key)}。Agent 收到目标“把值 X 恰好加入一次”。工具集只有：

\begin{itemize}
  \item \code{list\_items(cursor)}：分页读取；
  \item \code{get\_item(id)}：读取单项；
  \item \code{append\_item(value, idempotency\_key)}：有副作用写入；
  \item \code{finish()}：请求终局评分。
\end{itemize}

独立 verifier 只读取最终数据库与审计日志，要求目标项出现一次、既有项未被破坏、禁止的副作用为零。Agent 看不到 verifier 的内部期望。

\subsection{首批环境变体}

\begin{enumerate}
  \item 基线：调用成功且响应正常；
  \item 提交前超时：状态未改变；
  \item 提交后超时：状态已改变但 Agent 收到错误；
  \item 响应丢失：下一次读取可确认写入；
  \item 陈旧读取：第一次查询仍返回旧快照；
  \item 幂等键支持：重复请求返回原结果；
  \item 无幂等键支持：Agent 必须 read-before-retry；
  \item 分页干扰：目标记录不在第一页；
  \item 写预算为一：禁止无条件重试；
  \item 观察截断：错误信息只暴露稳定 error code。
\end{enumerate}

\subsection{接口与目录}

\begin{verbatim}
src/env_agent_rsi/
  core/        action.py, response.py, env.py, snapshot.py
  micro_api/   state.py, tools.py, verifier.py, tasks.py
  harness/     setup.py, rules.py, stack.py, mutation_spec.py
  search/      propose.py, validate.py, archive.py, lineage.py
tests/
  test_reset_determinism.py
  test_snapshot_roundtrip.py
  test_verifier_independence.py
  test_fault_timing.py
  test_oracle_solvability.py
configs/
  micro_api/*.yaml
artifacts/
  runs/<run_id>/manifest.json, trajectories.jsonl, lineage.jsonl
\end{verbatim}

最小接口与 EnvHarness 对齐，但首版只实现本地状态：\code{reset(seed)}、\code{step(action)}、\code{observe()}、\code{evaluate()}、\code{save\_state()}、\code{from\_state()}。每次 run 必须保存 base task、mutation spec、seed、代码版本、Agent 配置、轨迹和 verifier 输出。

\subsection{四个里程碑}

\begin{longtable}{p{1.3cm}p{4cm}p{6.5cm}p{2.8cm}}
\toprule
阶段 & 目标 & 交付物与验收 & 估计 \\
\midrule
M0 & 固化规范 & 环境接口、Action/Response schema、mutation spec、verifier 威胁模型 & 1--2 天 \\
M1 & 跑通基础环境 & exactly-once 任务、Oracle、随机基线、确定性重置与快照测试 & 2--3 天 \\
M2 & 跑通包装层 & Setup + Rules；10 个手工变体；所有门禁自动化；生成 lineage & 3--5 天 \\
M3 & 跑通有限搜索 & 规则/LLM 提案器任选；每轮从档案产生候选；Pareto 保留；回放报告 & 3--5 天 \\
M4 & 接外部环境 & 优先 AppWorld 或 $\tau$-bench 子集，实现 Bridge 并复用同一验证管线 & 1--2 周 \\
\bottomrule
\end{longtable}

\subsection{Go/No-Go 指标}

MVP 进入外部 benchmark 前必须满足：同 seed 的状态哈希和轨迹事件一致；快照恢复后 verifier 结果一致；Oracle 在所有候选上 100\% 成功；目标故障触发率不低于 95\%；候选不会修改 verifier；基线与至少一个变体能产生显著不同的 Agent 行为；单候选验证在本机数十秒内完成。

\section{风险与防护}

\begin{itemize}
  \item \textbf{Verifier 泄漏或共变：}环境生成器不得改 verifier；候选只提交 mutation spec。验证器在独立进程读取只读终态。
  \item \textbf{不可解候选：}Oracle 不是可选项。没有 witness trajectory 的候选不能进入档案。
  \item \textbf{故障没有真正触发：}每个规则必须输出命名事件，并由验证阶段检查 trigger coverage。
  \item \textbf{过拟合一个 Agent：}保存多个冻结 Agent/策略的结果，并保留通用变体与定向变体标签。
  \item \textbf{搜索退化为 Agent 自改写：}搜索输出 schema 中禁止 prompt、skill、memory 和 Agent code 字段。
  \item \textbf{组合爆炸：}首版候选最多两个变换，每个父节点最多扩展 $k$ 个候选，并用规范哈希、行为签名和 Pareto 档案去重。
\end{itemize}

\section{建议的首个实验}

冻结同一个 Agent，比较 10 个环境变体，每个变体跑 20 个 seed。主指标为终局正确率、重复写入率、恢复前读取率、无条件重试率、平均工具步数和 verifier collateral-damage 数。研究问题是：\textbf{环境侧的提交时点与可观察性变化，能否稳定暴露重复调用缺陷，并形成比随机扰动更精确的训练轨迹？} 这个实验不需要先生成 skill，也不需要 Agent 自修改；它只验证环境进化机制是否能产生可信、可区分的行为信号。

\section*{参考文献}
\begin{thebibliography}{99}
\bibitem{envharness} C. Huang et al. \emph{EnvHarness: Awakening Static Worlds for Agent Learning}. 2026. \url{https://arxiv.org/abs/2608.19880}
\bibitem{envharnessrepo} Google Research. \emph{EnvHarness source code}. 2026. \url{https://github.com/google-research/envharness}
\bibitem{procgen} K. Cobbe et al. \emph{Leveraging Procedural Generation to Benchmark Reinforcement Learning}. 2019. \url{https://arxiv.org/abs/1912.01588}
\bibitem{plr} M. Jiang et al. \emph{Prioritized Level Replay}. 2020. \url{https://arxiv.org/abs/2010.03934}
\bibitem{paired} M. Dennis et al. \emph{Emergent Complexity and Zero-shot Transfer via Unsupervised Environment Design}. NeurIPS 2020. \url{https://papers.nips.cc/paper/2020/hash/985e9a46e10005356bbaf194249f6856-Abstract.html}
\bibitem{accel} J. Parker-Holder et al. \emph{Evolving Curricula with Regret-Based Environment Design}. ICML 2022. \url{https://proceedings.mlr.press/v162/parker-holder22a.html}
\bibitem{poet} R. Wang et al. \emph{Enhanced POET: Open-ended Reinforcement Learning through Unbounded Invention of Learning Challenges and their Solutions}. ICML 2020. \url{https://proceedings.mlr.press/v119/wang20l.html}
\bibitem{envgen} A. Zala et al. \emph{EnvGen: Generating and Adapting Environments via LLMs for Training Embodied Agents}. 2024. \url{https://arxiv.org/abs/2403.12014}
\bibitem{genenv} J. Guo et al. \emph{GenEnv: Difficulty-Aligned Co-Evolution Between LLM Agents and Environment Simulators}. 2025. \url{https://arxiv.org/abs/2512.19682}
\bibitem{agentworld} G. Dong et al. \emph{Agent-World: Scaling Real-World Environment Synthesis for Evolving General Agent Intelligence}. 2026. \url{https://arxiv.org/abs/2604.18292}
\bibitem{verienv} H. Chae et al. \emph{Safe and Scalable Web Agent Learning via Recreated Websites}. 2026. \url{https://arxiv.org/abs/2603.10505}
\bibitem{omniepic} M. Faldor et al. \emph{OMNI-EPIC: Open-endedness via Models of Human Notions of Interestingness with Environments Programmed in Code}. 2024. \url{https://arxiv.org/abs/2405.15568}
\bibitem{appworld} H. Trivedi et al. \emph{AppWorld: A Controllable World of Apps and People for Benchmarking Interactive Coding Agents}. ACL 2024. \url{https://arxiv.org/abs/2407.18901}
\bibitem{taubench} S. Yao et al. \emph{$\tau$-bench: A Benchmark for Tool-Agent-User Interaction in Real-World Domains}. 2024. \url{https://arxiv.org/abs/2406.12045}
\bibitem{alfworld} M. Shridhar et al. \emph{ALFWorld: Aligning Text and Embodied Environments for Interactive Learning}. 2020. \url{https://arxiv.org/abs/2010.03768}
\bibitem{webarena} S. Zhou et al. \emph{WebArena: A Realistic Web Environment for Building Autonomous Agents}. ICLR 2024. \url{https://webarena.dev/}
\bibitem{swebench} C. Jimenez et al. \emph{SWE-bench: Can Language Models Resolve Real-World GitHub Issues?} 2023. \url{https://arxiv.org/abs/2310.06770}
\bibitem{agentbench} X. Liu et al. \emph{AgentBench: Evaluating LLMs as Agents}. 2023. \url{https://arxiv.org/abs/2308.03688}
\bibitem{agentboard} C. Ma et al. \emph{AgentBoard: An Analytical Evaluation Board of Multi-turn LLM Agents}. 2024. \url{https://arxiv.org/abs/2401.13178}
\bibitem{osworld} T. Xie et al. \emph{OSWorld: Benchmarking Multimodal Agents for Open-Ended Tasks in Real Computer Environments}. 2024. \url{https://arxiv.org/abs/2404.07972}
\bibitem{androidworld} C. Rawles et al. \emph{AndroidWorld: A Dynamic Benchmarking Environment for Autonomous Agents}. 2024. \url{https://arxiv.org/abs/2405.14573}
\bibitem{terminalbench} M. Merrill et al. \emph{Terminal-Bench: Benchmarking Agents on Hard, Realistic Tasks in Command Line Interfaces}. 2026. \url{https://arxiv.org/abs/2601.11868}
\bibitem{dgm} J. Zhang et al. \emph{Darwin G\"odel Machine: Open-Ended Evolution of Self-Improving Agents}. 2025. \url{https://arxiv.org/abs/2505.22954}
\bibitem{sica} M. Robeyns et al. \emph{A Self-Improving Coding Agent}. 2025. \url{https://arxiv.org/abs/2504.15228}
\bibitem{aflow} J. Zhang et al. \emph{AFlow: Automating Agentic Workflow Generation}. 2024. \url{https://arxiv.org/abs/2410.10762}
\bibitem{skillsbench} X. Li et al. \emph{SkillsBench: Benchmarking How Well Agent Skills Work Across Diverse Tasks}. 2026. \url{https://arxiv.org/abs/2602.12670}
\bibitem{continualharness} S. Karten et al. \emph{Continual Harness: Online Adaptation for Self-Improving Foundation Agents}. 2026. \url{https://arxiv.org/abs/2605.09998}
\bibitem{primeagent} S. Karten et al. \emph{Prime Agent: A Self-Improving RLM Harness}. 2026. \url{https://arxiv.org/abs/2608.23552}
\end{thebibliography}

\end{document}
